\scanpage{28}

$G$ opère sur $X$, \uline{stabilisateur} d'un $x\in X(k)\unsure{=|\simeq}\mathrm{Hom}_{\mathit{Aff}_k}(e_k,X)$

\[
g\longmapsto g.x
\]

\diagram{\begin{tikzcd}[ampersand replacement=\&,row sep=small]
G\arrow[r]\&X\\
\deleted{\unsure{K}}G_{(x)}\arrow[u]\arrow[r]\&e_k\arrow[u]
\end{tikzcd}}

$\mathrm I$.

Si on veut le sous-groupe de stabilité de plusieurs éléments $x_i\in X_i(k)$, $G$ opérant sur les $X_i$, on prend $\bigcap_i G_{x_i}$ (cas particulier de produit fibré) \unsure{où|ou} on fait opérer $G$ sur $X=\prod X_i$, et on prend le stabilisateur de $x=(x_i)_{i\in I}\in X(k)$.

\textcircled{10} $\varinjlim$ dans $\mathit{Aff}_k$. Ne s'interprète pas de façon commode en terme de $\widehat{\mathit{Aff}_k}\simeq\underline{\mathrm{Hom}}(\deleted{\unsure{k}}\mathit{Alg}_k,(\mathrm{Ens}))$, et n'a pas beaucoup de sens géométriquement, sauf en des cas très particuliers.

[\uline{Exemple} $G=\mathrm{Gl}(n)_k$, $B=\mathrm{Tr}(n)_k\subset G$, alors \deleted{\unsure{B}} $B$ définit une relation d'équivalence dans $G$, et le quotient $G/B$ \uline{dans} $\mathit{Aff}_k$ est $e_k$ -- ce qui est idiot !]

% slanted
\margin{définir la relation d'équivalence associée pour un \unsure{ss-groupe|sous-groupe} d'un \unsure{gr. alg.|groupe} \ldots}

Correspond à $\varprojlim$ dans $\mathit{Alg}_k$. Ce sont les $\varprojlim$ ensemblistes (nous fait une belle jambe !)

\deleted{\unsure{a)|$*$}} \uline{Sommes dans} $\mathit{Aff}_k$. Correspond aux produits dans $\mathit{Alg}_k$. Seules les sommes finies ($\Longleftrightarrow$ produits finis) ont une interprétation \uline{géométrique simple}.

\[
X_i\longrightarrow X
\]

Attention, $\boxed{\coprod_i X_i(k')\to X(k')}$ n'est pas bijective pour $\forall\ k'$ ; \deleted{\unsure{f}} i.e.

\[
\boxed{\coprod_{\unsure{I|i}}\mathrm{Hom}_{k\text{-}\mathrm{alg}}(\deleted{\unsure{B}}A_{\unsure{i|j}},k')\longrightarrow
\mathrm{Hom}_{k\text{-}\mathrm{alg}}(A=\prod A_i,k')}
\]

n'est pas bijective. \inserted{Elle est \uline{injective} si $k'\ne0$, / \deleted{\unsure{elle n'est pas}}} Un hom de $A$ dans $k'$ ``consiste'' en la donnée d'une famille d'idempotents de $k'$ $(e'_i)_{i\in I}$ de somme $1$ et deux à deux orthogonaux et d'hom. $A_i\to k'/(1-e'_i)k'$ \unsure{?}\ldots{} Mais si

% archive sheet number: 27
